% Wyniki skalarne przejęte z AM I-10 z oryginalnymi etykietami.
% Kryterium sum skończonych rozwija AF-27, lem:net-unconditional27.
\chapter{Zbieżność bezwzględna, warunkowa i bezwarunkowa}
\label{ma2:ch:unconditional-series}

Niech $E$ będzie przestrzenią unormowaną nad $\K=\R$ lub
$\K=\C$. Przez $\N_+$ rozumiemy zbiór dodatnich liczb
naturalnych.
Zbieżność szeregu zawsze oznacza zbieżność jego sum
częściowych w normie. Korzystamy z rozdziału
\ref{ma2:ch:banach-series} oraz teorii skalarnej z tomu I.

\section{Trzy rodzaje zbieżności}
\begin{definicja}
\label{ma2:def:unconditional-series}
Szereg $\sum_{n\ge1}x_n$ nazywamy zbieżnym warunkowo
w sensie zbieżności niebezwzględnej, gdy jest zbieżny,
ale $\sum_{n\ge1}\|x_n\|=+\infty$.
Przestawieniem szeregu nazywamy szereg
$\sum_{k\ge1}x_{\pi(k)}$, gdzie $\pi:\N_+\to\N_+$
jest bijekcją.
Szereg nazywamy zbieżnym bezwarunkowo do $S\in E$,
gdy dla każdej takiej bijekcji szereg przestawiony
jest zbieżny do $S$.
\end{definicja}
W literaturze spotyka się również określenie zbieżności
warunkowej dla szeregu zbieżnego, ale nie bezwarunkowo.
Aby uniknąć niejednoznaczności, będziemy dalej mówić
wprost o szeregu zbieżnym niebezwzględnie albo
zbieżnym, lecz nie bezwarunkowo. W skończonym wymiarze
te dwie własności są równoważne, co wykażemy w
twierdzeniu \ref{ma2:thm:finite-unconditional-absolute}.
W nieskończonym wymiarze mogą się różnić.

\begin{twierdzenie}[Przestawienia szeregów bezwzględnie zbieżnych]
\label{ma2:thm:absolute-rearrangement}
Jeżeli $E$ jest przestrzenią Banacha, a szereg
$\sum x_n$ jest bezwzględnie zbieżny, to jest bezwarunkowo
zbieżny. Każde przestawienie jest również bezwzględnie
zbieżne, a suma norm jest niezależna od przestawienia.
W dowolnej przestrzeni unormowanej ten sam wniosek
zachodzi, jeśli oprócz zbieżności szeregu norm założymy
istnienie sumy $\sum x_n$ w $E$.
\end{twierdzenie}
\begin{proof}
Niech $A=\sum_{n\ge1}\|x_n\|$ oraz $S=\sum_{n\ge1}x_n$;
istnienie $S$ wynika z twierdzenia \ref{ma2:thm:absolute-series}
w przypadku przestrzeni Banacha, a w pozostałym przypadku
jest założeniem. Oszacowanie
$\|S-\sum_{n=1}^N x_n\|\le\sum_{n>N}\|x_n\|$
w obu przypadkach wynika z nierówności trójkąta dla
sum skończonych i przejścia do granicy, jak w dowodzie
twierdzenia \ref{ma2:thm:absolute-series}.
Dla bijekcji $\pi$ niech $F_k=\{\pi(1),\ldots,\pi(k)\}$.
Z iniektywności i nieujemności norm mamy
$\sum_{j=1}^k\|x_{\pi(j)}\|\le
\sum_{n=1}^{\max F_k}\|x_n\|\le A$.
Sumy te są niemalejące, więc mają granicę $A'\le A$.
Dla każdego $N\ge1$ istnieje
$k_0=\max_{1\le n\le N}\pi^{-1}(n)$.
Dla $k\ge k_0$ zachodzi $\{1,\ldots,N\}\subseteq F_k$,
więc $A'\ge\sum_{n=1}^N\|x_n\|$.
Przechodząc do granicy $N\to\infty$, dostajemy $A'\ge A$,
czyli $A'=A$.
Dla $\varepsilon>0$ wybieramy $N$ tak, by
$\sum_{n>N}\|x_n\|<\varepsilon/2$.
Dla $k\ge k_0$ nierówność trójkąta oraz powyższe
oszacowanie reszty dają
\[
 \left\|\sum_{j=1}^k x_{\pi(j)}-S\right\|
 \le\left\|\sum_{n\in F_k\setminus\{1,\ldots,N\}}x_n\right\|
 +\left\|\sum_{n=1}^N x_n-S\right\|
 \le 2\sum_{n>N}\|x_n\|<\varepsilon.
\]
Zatem wszystkie przestawienia mają sumę $S$.
\end{proof}

\section{Szeregi naprzemienne}
\begin{twierdzenie}[Kryterium Leibniza]
\label{ma1:thm:leibniz-series}
Niech $(b_n)_{n\ge1}$ będzie nierosnącym ciągiem nieujemnych liczb rzeczywistych takim, że $b_n\to0$. Wówczas szereg
\[
 \sum_{n=1}^\infty (-1)^{n-1}b_n
\]
jest zbieżny. Jeżeli $S$ oznacza jego sumę, $S_0=0$, a $S_N=\sum_{n=1}^N(-1)^{n-1}b_n$ dla $N\ge1$, to dla każdego całkowitego $N\ge0$ zachodzi
\[
 0\le(-1)^N(S-S_N)\le b_{N+1}.
\]
W szczególności $|S-S_N|\le b_{N+1}$.
\end{twierdzenie}
\begin{proof}
Dla $k\ge0$ mamy $S_{2k+2}-S_{2k}=b_{2k+1}-b_{2k+2}\ge0$, więc ciąg parzystych sum częściowych jest niemalejący. Dla $k\ge0$ mamy również $S_{2k+3}-S_{2k+1}=-b_{2k+2}+b_{2k+3}\le0$, czyli ciąg nieparzystych sum częściowych jest nierosnący. Ponadto $S_{2k+1}-S_{2k}=b_{2k+1}\ge0$. Dla każdego $k\ge0$ otrzymujemy
\[
 0=S_0\le S_{2k}\le S_{2k+1}\le S_1=b_1.
\]
Z twierdzenia o zbieżności ciągów monotonicznych (tom I, rozdział o granicach ciągów) istnieją skończone granice $E=\lim_k S_{2k}$ i $O=\lim_k S_{2k+1}$. Ponieważ $2k+1\to\infty$, z założenia $b_n\to0$ otrzymujemy $S_{2k+1}-S_{2k}=b_{2k+1}\to0$, rachunek granic daje $O-E=0$. Oznaczmy wspólną granicę przez $S$. Dla dowolnego $\varepsilon>0$ istnieje całkowity $K\ge0$ taki, że dla każdego $k\ge K$ zarówno $|S_{2k}-S|<\varepsilon$, jak i $|S_{2k+1}-S|<\varepsilon$. Każdy całkowity indeks $N\ge2K$ ma postać $2k$ lub $2k+1$ z $k\ge K$. Zatem $S_N\to S$.

Monotoniczność i twierdzenie o granicach ciągów monotonicznych dają $S=\sup_{k\ge0}S_{2k}=\inf_{k\ge0}S_{2k+1}$, a więc $S_{2k}\le S\le S_{2k+1}$ dla każdego $k\ge0$. Dla parzystego $N=2k$ wynika stąd $0\le S-S_N\le S_{2k+1}-S_{2k}=b_{N+1}$. Dla nieparzystego $N=2k+1$ korzystamy z $S_{2k+2}\le S\le S_{2k+1}$, więc $0\le S_N-S\le S_{2k+1}-S_{2k+2}=b_{N+1}$. Oba przypadki dają żądane oszacowanie.
\end{proof}
\begin{przyklad}
\label{ma1:ex:alternating-harmonic}
Szereg $\sum_{n=1}^\infty(-1)^{n-1}/n$ jest zbieżny, ponieważ $(1/n)$ jest nierosnący i zbiega do zera. Nie jest zbieżny bezwzględnie: szereg modułów jest szeregiem harmonicznym (tom I, rozdział 09). Dla jego sumy $S$ mamy
$S_4=7/12\le S\le S_5=47/60$.
Dla błędu mniejszego od $1/100$ wystarczy obliczyć $S_{100}$, ponieważ $|S-S_{100}|\le1/101<1/100$. Nie potrzebujemy tu identyfikacji $S$ z $\log 2$.
\end{przyklad}

\begin{figure}[htbp]
\centering
\begin{tikzpicture}[x=.75cm,y=4.2cm,>=Stealth]
\draw[->] (-.15,0)--(10.9,0) node[right] {$N$};
\draw[->] (0,-.04)--(0,1.13) node[above] {$S_N$};
\foreach \n in {1,...,10} {
 \draw (\n,.012)--(\n,-.012) node[below] {\small $\n$};
}
\draw[dashed,gray] (0,1)--(10.5,1);
\node[left] at (0,1) {$1$};
\draw[blue!70!black,thick] (0,0)--(2,.5)--(4,.583333)--(6,.616667)--(8,.634524)--(10,.645635);
\draw[orange!85!black,thick] (1,1)--(3,.833333)--(5,.783333)--(7,.759524)--(9,.745635);
\foreach \n/\s in {0/0,2/.5,4/.583333,6/.616667,8/.634524,10/.645635}
 \fill[blue!70!black] (\n,\s) circle (1.8pt);
\foreach \n/\s in {1/1,3/.833333,5/.783333,7/.759524,9/.745635}
 \fill[orange!85!black] (\n,\s) circle (1.8pt);
\node[blue!70!black,below] at (8,.60) {$S_{2k}$};
\node[orange!85!black,above] at (8,.81) {$S_{2k+1}$};
\draw[densely dotted,gray] (9,.745635)--(10.45,.745635);
\draw[densely dotted,gray] (10,.645635)--(10.45,.645635);
\draw[<->,gray] (10.35,.645635)--(10.35,.745635);
\node[right] at (10.45,.695635) {\small $1/10$};
\end{tikzpicture}
\caption{Parzyste i nieparzyste sumy częściowe szeregu naprzemiennego harmonicznego. Odcinki łączą punkty wyłącznie dla czytelności. Długość przedziału między $S_{10}$ a $S_9$ wynosi $1/10$.}
\label{ma1:fig:leibniz-sums}
\end{figure}
\FloatBarrier

\section{Część dodatnia i ujemna szeregu rzeczywistego}
\begin{definicja}
\label{ma2:def:positive-negative-parts}
Dla liczby rzeczywistej $a$ określamy $a^+=\max(a,0)$ i $a^-=\max(-a,0)$. Obie liczby są nieujemne. Dla $a\ge0$ mamy $a^+=a$, $a^-=0$, a dla $a<0$ mamy $a^+=0$, $a^-=-a$. W obu przypadkach zachodzą $a=a^+-a^-$ oraz $|a|=a^++a^-$.
\end{definicja}
\begin{twierdzenie}
\label{ma1:thm:positive-negative-parts}
Niech $(a_n)$ będzie ciągiem rzeczywistym. Szereg $\sum a_n$ jest zbieżny bezwzględnie wtedy i tylko wtedy, gdy oba szeregi $\sum a_n^+$ i $\sum a_n^-$ są zbieżne. Jeżeli $\sum a_n$ jest zbieżny warunkowo, to obie sumy częściowe tych szeregów dążą do $+\infty$.
\end{twierdzenie}
\begin{proof}
Jeżeli $\sum|a_n|$ jest zbieżny, z $0\le a_n^\pm\le |a_n|$ oraz kryterium porównawczego wynika zbieżność obu szeregów. Odwrotnie, jeśli oba są zbieżne, ich suma $\sum(a_n^++a_n^-)=\sum|a_n|$ jest zbieżna z liniowości.

Załóżmy teraz zbieżność warunkową. Niech $P_N=\sum_{n=1}^N a_n^+$, $Q_N=\sum_{n=1}^N a_n^-$ oraz $S_N=\sum_{n=1}^N a_n$. Mamy $S_N=P_N-Q_N$. Gdyby $(P_N)$ miało skończoną granicę, to $Q_N=P_N-S_N$ również miałoby skończoną granicę, wbrew brakowi zbieżności bezwzględnej. Gdyby natomiast $(Q_N)$ miało skończoną granicę, to $P_N=S_N+Q_N$ również miałoby skończoną granicę, z tym samym skutkiem. Oba ciągi są niemalejące, więc na mocy kryterium zbieżności szeregów nieujemnych z tomu I, rozdział 09, dążą do $+\infty$.
\end{proof}

\section{Twierdzenie Riemanna o przestawieniach}
\begin{twierdzenie}[Riemanna]
\label{ma1:thm:riemann-rearrangement}
Niech $(a_n)_{n\ge1}$ będzie ciągiem rzeczywistym takim, że szereg $\sum_{n=1}^\infty a_n$ jest zbieżny warunkowo (czyli zbieżny, lecz nie bezwzględnie). Dla każdej liczby $s\in\R$ istnieje bijekcja $\pi:\N_+\to\N_+$, dla której $\sum_{k=1}^\infty a_{\pi(k)}=s$. Istnieją również przestawienia, których sumy częściowe dążą odpowiednio do $+\infty$ i do $-\infty$, oraz przestawienie, którego sumy częściowe nie mają granicy nawet w $\overline{\R}$.
\end{twierdzenie}
\begin{proof}
Z twierdzenia \ref{ma1:thm:positive-negative-parts} wynika, że zbiór indeksów wyrazów dodatnich i zbiór indeksów wyrazów ujemnych są nieskończone: gdyby którykolwiek był skończony, odpowiedni ciąg sum częściowych części dodatniej lub ujemnej byłby od pewnego indeksu stały, wbrew rozbieżności do $+\infty$. Wyliczmy je rosnąco jako $(p_j)_{j\ge1}$ i $(q_j)_{j\ge1}$. Piszemy $u_j=a_{p_j}>0$ oraz $v_j=-a_{q_j}>0$. Każda suma częściowa części dodatniej szeregu jest sumą pierwszych kilku $u_j$ (lub sumą pustą), a dla każdego $J\ge1$ zachodzi $\sum_{j=1}^J u_j=\sum_{n=1}^{p_J}a_n^+$. Ponieważ $p_J\ge J$, prawa strona dąży do $+\infty$. Analogicznie $\sum_{j=1}^J v_j=\sum_{n=1}^{q_J}a_n^-\to+\infty$. Piszemy więc $\sum_{j=1}^\infty u_j=\sum_{j=1}^\infty v_j=+\infty$ w sensie rozbieżności do nieskończoności. Ponadto $a_n\to0$ na mocy warunku koniecznego zbieżności. Ponieważ $p_j,q_j\ge j$, mamy $a_{p_j}\to0$ i $a_{q_j}\to0$ jako podciągi ciągu zbieżnego do zera. Rachunek granic daje więc $u_j,v_j\to0$. Zera wyliczymy osobno rosnąco według indeksów; ich wstawianie nie zmienia wartości sum częściowych.

Ustalmy najpierw $s\in\R$. Zaczynamy od sumy $0$. Jeżeli $0>s$, dodajemy po kolei nieużyte wyrazy $-v_j$ do pierwszej chwili, w której suma stanie się mniejsza od $s$. Taki skończony krok istnieje, ponieważ $\sum_{j=1}^J v_j\to+\infty$ przy $J\to\infty$. W przeciwnym razie zaczynamy bez tego wstępnego kroku. Następnie dodajemy co najmniej jeden kolejny nieużyty wyraz $u_j$, kończąc blok w pierwszej chwili, gdy suma przekroczy $s$. Od tej chwili dodajemy kolejne nieużyte wyrazy $-v_j$, kończąc blok w pierwszej chwili, gdy suma spadnie poniżej $s$. Powtarzamy oba kroki w nieskończoność. Każdy blok jest skończony: sumy pozostałych dodatnich wyrazów i modułów pozostałych ujemnych wyrazów są nieograniczone po pominięciu skończonego początku. Każdy z dwóch rodzajów bloków zawiera co najmniej jeden wyraz.

Oznaczmy przez $U_k$ sumę częściową na końcu $k$-tego dodatniego bloku, a przez $V_k$ sumę na końcu następującego po nim ujemnego bloku. Niech $i_k,j_k$ będą numerami ostatnich wykorzystanych wyrazów odpowiednio w ciągach $(u_j)$ i $(v_j)$. Bezpośrednio przed ostatnim dodaniem wyrazu dodatniego suma jest nie większa od $s$, a bezpośrednio przed ostatnim dodaniem wyrazu ujemnego jest nie mniejsza od $s$. Wynika to z wyboru pierwszego przekroczenia progu (suma początkowa pierwszego dodatniego bloku jest również nie większa od $s$). Otrzymujemy
\[
 0<U_k-s\le u_{i_k},\qquad 0<s-V_k\le v_{j_k}.
\]
Ponieważ każdy blok zawiera co najmniej jeden nowy wyraz danego znaku, $i_k,j_k\ge k$. Dla dowolnego $\varepsilon>0$ wybieramy $K$ tak, aby $u_j,v_j<\varepsilon$ dla $j\ge K$. Dla $k\ge K$ oba przekroczenia progu są mniejsze od $\varepsilon$, zatem $U_k\to s$ i $V_k\to s$. Wszystkie sumy częściowe wewnątrz kolejnych bloków dodatnich leżą pomiędzy $V_{k-1}$ a $U_k$ (dla pierwszego bloku używamy sumy po kroku wstępnym). Wewnątrz bloków ujemnych leżą pomiędzy $V_k$ a $U_k$. Dla dowolnego $\varepsilon>0$ wszystkie te krańce, poza skończenie wieloma początkowymi blokami, należą do $(s-\varepsilon,s+\varepsilon)$. Każdy blok jest skończony, więc również cały ciąg sum częściowych zbiega do $s$.

Ponieważ w każdym cyklu wykorzystujemy co najmniej jeden nowy wyraz każdego znaku, a wyrazy danego znaku wybieramy według kolejności ich indeksów, każdy indeks wyrazu niezerowego zostanie użyty dokładnie raz. Po każdym pełnym cyklu wstawiamy jeden kolejny nieużyty indeks wyrazu zerowego, jeżeli taki jeszcze istnieje. Każdy indeks zera również zostanie użyty dokładnie raz. Powstałe wyliczenie wszystkich dodatnich liczb naturalnych bez powtórzeń określa bijekcję $\pi$; wstawienie zer nie zmienia wniosku o granicy.

Dla granicy $+\infty$ w $k$-tym cyklu dodajemy najpierw co najmniej jeden kolejny dodatni wyraz, a następnie tyle kolejnych dodatnich wyrazów, ile potrzeba, aby suma przekroczyła $k+v_k$. Dodajemy wtedy dokładnie jeden ujemny wyraz $-v_k$. Po tym kroku suma jest większa od $k$. W każdym cyklu suma najpierw rośnie, a potem spada do wartości większej od $k$. Dla $k\ge2$ wszystkie sumy wewnątrz cyklu są więc większe od $k-1$, bo początkowa suma jest większa od $k-1$. Po każdym cyklu można ponownie wstawić jedno zero. Wykorzystujemy wszystkie indeksy każdego znaku i każdy indeks zera, jak wyżej. Dla dowolnego $M\in\R$ wybieramy całkowite $K\ge2$ takie, że $K>M+1$; wszystkie sumy po pierwszych $K-1$ cyklach są większe od $M$. Otrzymujemy granicę $+\infty$. Przez zastosowanie konstrukcji do $(-a_n)$ uzyskujemy przestawienie o granicy $-\infty$.

Wreszcie możemy w $k$-tym cyklu dodawać kolejno wyrazy dodatnie do pierwszego przekroczenia poziomu $k$, a następnie wyrazy ujemne do pierwszego spadku poniżej $-k$. Oba bloki są skończone z rozbieżności odpowiednich pozostałych sum. Są niepuste, ponieważ pierwszy cykl rozpoczynamy od $0<1$, każdy następny od sumy mniejszej od $-(k-1)<k$, a każdy blok ujemny od sumy większej od $k>-k$. Dodajemy zera według tej samej reguły. Znowu otrzymujemy bijekcję, ale sumy częściowe na końcach dodatnich bloków dążą do $+\infty$, a na końcach ujemnych do $-\infty$. Granice dolna i górna tego ciągu sum wynoszą odpowiednio $-\infty$ i $+\infty$, ponieważ istnieją podciągi dążące do obu tych wartości. Zatem ciąg nie jest zbieżny w $\overline{\R}$.
\end{proof}

\begin{przyklad}
\label{ma1:ex:harmonic-rearrangement}
Niech $S=\sum_{n=1}^\infty(-1)^{n-1}/n$. Przestawmy wyrazy szeregu, umieszczając po dwa kolejne dodatnie wyrazy, a następnie jeden kolejny ujemny wyraz:
\[
 1+\frac13-\frac12+\frac15+\frac17-\frac14+\cdots.
\]
W $k$-tym bloku występują wyrazy o indeksach $4k-3,4k-1,2k$ szeregu wyjściowego. Każdy nieparzysty indeks ma dokładnie jedną z postaci $4k-3$ albo $4k-1$ dla $k\ge1$, a każdy parzysty indeks ma jednoznacznie postać $2k$ dla $k\ge1$. Zbiory tych trzech rodzajów indeksów są parami rozłączne, więc jest to przestawienie.
Oznaczmy sumy częściowe szeregu wyjściowego przez $S_N$, a przestawionego przez $T_N$. Dla $m\ge1$ zachodzi
\[
 T_{3m}-S_{4m}
 =\sum_{k=m+1}^{2m}\frac1{2k}
 =\frac12\sum_{k=m+1}^{2m}\frac1k.
\]
Istotnie, w obu sumach występuje dokładnie $2m$ pierwszych wyrazów dodatnich, natomiast $T_{3m}$ zawiera tylko $m$ pierwszych wyrazów ujemnych, a $S_{4m}$ zawiera ich $2m$. Różnica jest więc sumą modułów brakujących wyrazów ujemnych. Nie obliczamy tej granicy za pomocą logarytmu. Zauważmy natomiast, że
\[
 S_{2m}=\sum_{k=1}^{2m}\frac1k
 -2\sum_{k=1}^m\frac1{2k}
 =\sum_{k=m+1}^{2m}\frac1k.
\]
Stąd $T_{3m}=S_{4m}+\tfrac12 S_{2m}\to\tfrac32 S$. Ponieważ moduły wyrazów w kolejnych blokach dążą do zera, długość bloku jest stała i wynosi $3$, więc wniosek \ref{ma2:cor:bounded-blocks} daje $T_N\to\tfrac32 S$. Z kryterium Leibniza mamy $S\ge S_2=1/2$, zatem $\tfrac32 S\ne S$.
\end{przyklad}


\section{Kryterium sum po zbiorach skończonych}
Dla skończonego zbioru $F\subseteq\N_+$ oznaczmy
$S_F=\sum_{n\in F}x_n$, przy czym $S_\varnothing=0$.
Suma skończona nie zależy od kolejności składników.
Piszemy $S_F\to S$ względem inkluzji, jeśli dla każdego
$\varepsilon>0$ istnieje skończony $F_0$ taki, że
$\|S_F-S\|<\varepsilon$ dla wszystkich skończonych
$F\supseteq F_0$. Jest to zbieżność sieci sum skończonych.

\begin{twierdzenie}[Kryterium bezwarunkowej zbieżności]
\label{ma2:thm:unconditional-finite}
W przestrzeni Banacha następujące warunki są równoważne:
\begin{enumerate}[label=(\roman*)]
\item każdy przestawiony szereg jest zbieżny (bez
zakładania równości jego sum z sumą pierwotną);
\item szereg jest bezwarunkowo zbieżny;
\item istnieje $S\in E$, dla którego $S_F\to S$
względem inkluzji;
\item dla każdego $\varepsilon>0$ istnieje całkowite
$N\ge1$ takie, że dla każdego skończonego
$F\subseteq\{N+1,N+2,\ldots\}$ zachodzi
$\|S_F\|<\varepsilon$;
\item każdy podszereg (wybrany przez ściśle rosnący
ciąg dodatnich indeksów) jest zbieżny.
\end{enumerate}
Sumy z (ii) i (iii) są równe sumie szeregu wyjściowego.
\end{twierdzenie}
\begin{proof}
Implikacja (ii)$\Rightarrow$(i) wynika z definicji.

(i)$\Rightarrow$(iv): załóżmy, że (iv) jest fałszywe.
Istnieje $\varepsilon>0$ takie, że dla każdego $N\ge1$
można znaleźć skończony zbiór $F\subseteq\{N+1,\ldots\}$
z $\|S_F\|\ge\varepsilon$. Budujemy przestawienie etapami.
Przed etapem $k$ zbiór $D$ wszystkich wypisanych indeksów
jest skończony. Wybieramy $N=\max(D\cup\{k\})$,
a następnie odpowiedni niepusty zbiór $F_k$ poza tym
początkiem. Wypisujemy wszystkie jego elementy rosnąco,
po czym wypisujemy najmniejszy jeszcze niewypisany indeks.
Każdy etap jest skończony i żaden indeks nie powtarza się.
Po $k$ etapach wszystkie indeksy od $1$ do $k$ są wypisane:
indukcyjnie indeks $k$ albo już został wypisany, albo
jest najmniejszym brakującym indeksem i zostaje dopisany.
Powstaje więc bijekcja $\pi:\N_+\to\N_+$.
Sumy częściowe tuż przed i po bloku $F_k$ różnią się
o $S_{F_k}$, którego norma jest co najmniej $\varepsilon$.
Ponieważ oba odpowiadające im indeksy sum częściowych
dążą do nieskończoności, warunek Cauchy'ego nie zachodzi.
Przeczy to (i).

(iv)$\Rightarrow$(iii): dla $\varepsilon>0$ dobieramy $N$
z (iv). Jeśli $q\ge p\ge N+1$, to
$\|\sum_{n=p}^q x_n\|<\varepsilon$, ponieważ
$\{p,\ldots,q\}\subseteq\{N+1,N+2,\ldots\}$. Kryterium Cauchy'ego
(twierdzenie \ref{ma2:thm:series-cauchy}) i zupełność
dają istnienie $S=\lim_M S_{\{1,\ldots,M\}}$.
Teraz dobieramy $N$ tak, aby dla wszystkich skończonych
$F\subseteq\{N+1,\ldots\}$ było $\|S_F\|<\varepsilon/2$.
Dla każdego $M>N$ mamy
$\|S_{\{1,\ldots,M\}}-S_{\{1,\ldots,N\}}\|<\varepsilon/2$.
Przejście do granicy daje
$\|S-S_{\{1,\ldots,N\}}\|\le\varepsilon/2$.
Dla skończonego $G\supseteq\{1,\ldots,N\}$ otrzymujemy
\[
 \|S_G-S\|\le\|S_{G\setminus\{1,\ldots,N\}}\|
   +\|S_{\{1,\ldots,N\}}-S\|<\varepsilon.
\]
To jest (iii). Nie potrzebujemy ogólnego twierdzenia
o zupełności sieci, a jedynie zupełności ciągowej.

(iii)$\Rightarrow$(ii): niech $\pi$ będzie bijekcją.
Dla $\varepsilon>0$ wybieramy skończony $F_0$ z (iii).
Jeśli $F_0\ne\varnothing$, to dla
$k\ge\max_{n\in F_0}\pi^{-1}(n)$ zbiór
$\{\pi(1),\ldots,\pi(k)\}$ zawiera $F_0$; jeśli
$F_0=\varnothing$, inkluzja zachodzi dla każdego $k$.
Oznacza to, że odpowiednia suma częściowa jest odległa
od $S$ o mniej niż $\varepsilon$.

(iv)$\Rightarrow$(v): dla rosnących indeksów $n_j$
mamy $n_j\ge j$. Po wyborze $N$ z (iv), każda suma
$\sum_{j=p}^q x_{n_j}$ przy $q\ge p>N$ jest sumą po
skończonym zbiorze poza pierwszymi $N$ indeksami.
Zatem kryterium Cauchy'ego i zupełność dają zbieżność.

(v)$\Rightarrow$(iv): jeśli (iv) jest fałszywe, dla
pewnego $\varepsilon>0$ wybieramy rekurencyjnie niepuste
skończone zbiory $F_k$ z $\|S_{F_k}\|\ge\varepsilon$
tak, aby każdy indeks w $F_{k+1}$ był większy od
$\max F_k$. Wystarczy zastosować negację (iv) z
$N=\max F_k$ (w pierwszym kroku z $N=1$).
Wyliczenie elementów kolejnych bloków w porządku
rosnącym daje ściśle rosnący ciąg indeksów.
Jeśli $m_k=\sum_{j=1}^k |F_j|$ i $m_0=0$, to $m_k\ge k$
i różnica sum częściowych podszeregu o indeksach
$m_k$ i $m_{k-1}$ wynosi $S_{F_k}$. Dla dowolnego
proponowanego progu warunku Cauchy'ego wybieramy $k$
tak duże, by $m_{k-1}$ przekraczało próg. Norma różnicy
jest co najmniej $\varepsilon$, więc warunek Cauchy'ego
nie zachodzi. Sprzeczność dowodzi (iv).

W (iii) sumy po $\{1,\ldots,N\}$ zbiegają do $S$,
bo każdy skończony $F_0$ jest zawarty w takim zbiorze
dla dostatecznie dużego $N$. W (ii) wystarczy wziąć
bijekcję identycznościową.
\end{proof}

\begin{uwaga}
Zbieżność wszystkich przestawień implikuje ich wspólną
sumę również w niezupełnej przestrzeni unormowanej:
dowód (i)$\Rightarrow$(iv) nie korzysta z zupełności,
a w dowodzie (iv)$\Rightarrow$(iii) granica zwykłych
sum częściowych istnieje wtedy z (i). Sam warunek
(iv) bez założenia zupełności nie gwarantuje istnienia
sumy w $E$.
\end{uwaga}

\begin{twierdzenie}[Ograniczone mnożniki]
\label{ma2:thm:unconditional-multipliers}
Jeżeli $\sum x_n$ jest bezwarunkowo zbieżny w przestrzeni
Banacha $E$, a $(a_n)$ jest ograniczonym ciągiem skalarów,
to $\sum a_nx_n$ jest bezwarunkowo zbieżny.
\end{twierdzenie}
\begin{proof}
Najpierw dla skończonego zbioru $F$ i rzeczywistych
współczynników $0\le t_n\le1$ wykażemy oszacowanie
\[
 \left\|\sum_{n\in F}t_n x_n\right\|
 \le \max_{G\subseteq F}\|S_G\|.
\]
Dla $F=\varnothing$ obie strony są zerowe. Dla
$F=\{n_1,\ldots,n_m\}$ porządkujemy indeksy tak, by
$t_{n_1}\ge\cdots\ge t_{n_m}$. Dla wygody oznaczmy
$b_k=t_{n_k}$ dla $1\le k\le m$ i $b_{m+1}=0$.
Wówczas
\[
 \sum_{j=1}^m t_{n_j}x_{n_j}
 =\sum_{k=1}^m(b_k-b_{k+1})
               \sum_{j=1}^k x_{n_j}.
\]
Współczynnik przy $x_{n_j}$ po prawej to suma
różnic od $k=j$ do $m$, czyli $b_j=t_{n_j}$.
Wszystkie różnice są nieujemne, a ich suma wynosi
$t_{n_1}\le1$. Nierówność trójkąta daje oszacowanie.

Niech $M=\sup_n|a_n|$. Gdy $M=0$, wszystkie
wyrazy są zerowe. Dla $M>0$ każdy $a_n/M$
w przypadku rzeczywistym jest różnicą dwóch liczb
z $[0,1]$, gdyż $|a_n/M|\le1$. W przypadku zespolonym zapisujemy go jako
$(t_n^+-t_n^-)+\mathrm i(s_n^+-s_n^-)$,
gdzie $t_n^\pm,s_n^\pm\in[0,1]$ są częściami dodatnimi
i ujemnymi odpowiednich współrzędnych, których moduły
nie przekraczają $|a_n/M|\le1$.
Z powyższego oszacowania wynika w obu przypadkach
\[
 \left\|\sum_{n\in F}a_nx_n\right\|
 \le 4M\max_{G\subseteq F}\|S_G\|.
\]
Dla $\varepsilon>0$ stosujemy warunek (iv) twierdzenia
\ref{ma2:thm:unconditional-finite} z dokładnością
$\varepsilon/(4M)$. Jeśli $F$ leży poza wybranym
początkiem, każdy jego podzbiór również tam leży,
więc skończone maksimum jest mniejsze od
$\varepsilon/(4M)$. Otrzymujemy warunek (iv)
dla szeregu $\sum a_nx_n$, a zatem jego bezwarunkową
zbieżność.
\end{proof}

\section{Wymiar skończony i nieskończony}
\begin{twierdzenie}
\label{ma2:thm:finite-unconditional-absolute}
W skończenie wymiarowej przestrzeni unormowanej
zbieżność bezwarunkowa jest równoważna zbieżności
bezwzględnej.
\end{twierdzenie}
\begin{proof}
Przestrzeń $E$ jest zupełna z twierdzenia
\ref{ma2:thm:finite}, więc zbieżność bezwzględna
implikuje bezwarunkową na mocy twierdzenia
\ref{ma2:thm:absolute-rearrangement}.
Dla wymiaru $0$ teza zachodzi, ponieważ wszystkie
wyrazy są zerowe. Niech zatem $d=\dim E\ge1$
i niech $(e_j)_{j=1}^d$ będzie bazą.
Zapiszmy $x_n=\sum_{j=1}^d a_{nj}e_j$.
Jeśli $\sum x_n$ jest bezwarunkowo zbieżny, to
z twierdzenia \ref{ma2:thm:series-coordinates}
każdy szereg współrzędnych jest zbieżny przy
każdym przestawieniu.
Dla rzeczywistego szeregu skalarnego, który jest
zbieżny, lecz nie bezwzględnie, twierdzenie
\ref{ma1:thm:riemann-rearrangement} daje rozbieżne
przestawienie. Zatem w przypadku rzeczywistym
każdy szereg współrzędnych jest bezwzględnie zbieżny.
W przypadku zespolonym dla każdej współrzędnej
piszemy $a_{nj}=u_{nj}+\mathrm i v_{nj}$.
Ponieważ $|\operatorname{Re}z|,|\operatorname{Im}z|\le |z|$,
zbieżność dowolnego przestawienia szeregu zespolonego
implikuje zbieżność odpowiadających mu przestawień obu
szeregów rzeczywistych. Z poprzedniego argumentu
$\sum_n|u_{nj}|<\infty$ i $\sum_n|v_{nj}|<\infty$,
a z $|a_{nj}|\le|u_{nj}|+|v_{nj}|$ wynika
$\sum_n|a_{nj}|<\infty$. Dla każdego $N$ mamy
\[
 \sum_{n=1}^N\|x_n\|
 \le \sum_{j=1}^d\|e_j\|\sum_{n=1}^N|a_{nj}|
 \le \sum_{j=1}^d\|e_j\|\sum_{n=1}^\infty|a_{nj}|<\infty.
\]
W pierwszej nierówności zastosowaliśmy jednorodność
i nierówność trójkąta do każdego $x_n$, a następnie
zamieniliśmy kolejność dwóch sum skończonych.
Prawa strona jest niezależna od $N$, więc niemalejące
sumy częściowe szeregu norm są zbieżne.
\end{proof}
Twierdzenie Riemanna w pełnym zakresie dowolnej
zadanej sumy dotyczy szeregów rzeczywistych.
Nie wynika z niego możliwość otrzymania dowolnej
sumy zespolonej z warunkowo zbieżnego szeregu:
przykładowo wszystkie sumy częściowe przestawień
rzeczywistego szeregu leżą na osi rzeczywistej.

\begin{przyklad}[Bezwarunkowa, ale nie bezwzględna zbieżność]
\label{ma2:ex:unconditional-l2}
W przestrzeni Banacha $\ell^2$ niech $e_n$ oznacza
standardowy wektor jednostkowy i niech $x_n=e_n/n$.
Mamy $\|x_n\|_2=1/n$, więc szereg norm jest rozbieżny.
Natomiast dla skończonego $F\subseteq\{N+1,N+2,\ldots\}$
\[
 \left\|\sum_{n\in F}\frac{e_n}{n}\right\|_2^2
 =\sum_{n\in F}\frac1{n^2}\le\sum_{n>N}\frac1{n^2}.
\]
Ostatni ogon dąży do zera, ponieważ szereg
$\sum n^{-2}$ jest zbieżny (tom I, rozdział 09).
Dla $\varepsilon>0$ wybieramy $N$ z ogonem
mniejszym niż $\varepsilon^2$, co daje (iv)
z twierdzenia \ref{ma2:thm:unconditional-finite}.
Szereg jest więc bezwarunkowo zbieżny, lecz nie
bezwzględnie zbieżny.
Jego suma to $x=(1/j)_{j\ge1}\in\ell^2$:
dla zwykłych sum częściowych
$\|x-S_N\|_2^2=\sum_{j>N}1/j^2\to0$.
\end{przyklad}

\begin{uwaga}
Zbieżność bezwzględna implikuje bezwarunkową w przestrzeni
Banacha, a bezwarunkowa implikuje zwykłą. Pierwsza
implikacja nie odwraca się w ogólności, jak pokazuje
przykład \ref{ma2:ex:unconditional-l2}; druga nie odwraca
się nawet w $\R$, na mocy twierdzenia Riemanna i
przykładu \ref{ma1:ex:alternating-harmonic}.
Nie twierdzimy, że każde nieskończenie wymiarowe $E$
ma taki kontrprzykład, gdyż ogólny wynik
Dvoretzky'ego--Rogersa wymaga osobnego dowodu.
\end{uwaga}
